\par\Needspace{6\baselineskip}\medskip
\begin{center}\hypertarget{native-231-article168}{}\pdfbookmark[2]{168.}{article-168}168.\end{center}
Verbindet man die Gleichungen 1 und 2 aus Artikel 112, so erhält man in der hier verwendeten Schreibweise die folgenden Gleichungen.
\[
\begin{aligned}
0={}&(n12)R\cos B\sin(l-a)\\
&-(n02)\bigl(\rho'\cos\beta'\sin(a'-a)+R'\cos B'\sin(l'-a)\bigr)\\
&+(n01)\bigl(\rho''\cos\beta''\sin(a''-a)+R''\cos B''\sin(l''-a)\bigr)
\end{aligned}
\]
\SourceLeaf{232}
\[
\begin{aligned}
0={}&(n23)\bigl(\rho'\cos\beta'\sin(a'''-a')+R'\cos B'\sin(a'''-l')\bigr)\\
&-(n13)\bigl(\rho''\cos\beta''\sin(a'''-a'')+R''\cos B''\sin(a'''-l'')\bigr)\\
&+(n12)R'''\cos B'''\sin(a'''-l''').
\end{aligned}
\]

Setzen wir
\[
\begin{gathered}
\frac{R'\cos B'\sin(l'-a)}{\cos\beta'\sin(a'-a)}-R'\cos\delta'=b'\\[.7em]
\frac{R\cos B\sin(l-a)}{\cos\beta''\sin(a''-a)}=\lambda\\[.7em]
\frac{R''\cos B''\sin(a'''-l'')}{\cos\beta''\sin(a'''-a'')}-R''\cos\delta''=b''\\[.7em]
\frac{R'''\cos B'''\sin(a'''-l''')}{\cos\beta'\sin(a'''-a')}=\lambda'''\\[.7em]
\frac{R'\cos B'\sin(a'''-l')}{\cos\beta'\sin(a'''-a')}-R'\cos\delta'=\varkappa'\\[.7em]
\frac{\cos\beta'\sin(a'-a)}{\cos\beta''\sin(a''-a)}=\mu'\\[.7em]
\frac{R''\cos B''\sin(l''-a)}{\cos\beta''\sin(a''-a)}-R''\cos\delta''=\varkappa''\\[.7em]
\frac{\cos\beta''\sin(a'''-a'')}{\cos\beta'\sin(a'''-a')}=\mu''
\end{gathered}
\]
und formen entsprechend um, so gehen diese Gleichungen über in
\[
\begin{gathered}
\frac{\mu'(1+P')(x'+b')}{1+\dfrac{Q'}{(x'x'+a'a')^{\frac32}}}=x''+\varkappa''+\lambda P'\\[1em]
\frac{\mu''(1+P'')(x''+b'')}{1+\dfrac{Q''}{(x''x''+a''a'')^{\frac32}}}=x'+\varkappa'+\lambda'''P'',
\end{gathered}
\]
oder, wenn wir außerdem
\[
\begin{gathered}
-\varkappa''-\lambda P'=c',\qquad\mu'(1+P')=d'\\[.5em]
-\varkappa'-\lambda'''P''=c'',\qquad\mu''(1+P'')=d'',
\end{gathered}
\]
setzen, in die folgenden
\[
\tag*{I.}x''=c'+\frac{d'(x'+b')}{1+\dfrac{Q'}{(x'x'+a'a')^{\frac32}}}
\]
\[
\tag*{II.}x'=c''+\frac{d''(x''+b'')}{1+\dfrac{Q''}{(x''x''+a''a'')^{\frac32}}}\mathpunct{\cdot}
\]
Mithilfe dieser beiden Gleichungen lassen sich $x'$ und $x''$ aus \newline $a',b',c',d',Q',a'',b'',c'',d'',Q''$ bestimmen. Müsste man daraus $x'$ oder $x''$ eliminieren, so gelangte man zu einer Gleichung sehr hohen Grades\HistoricalNoteEight{Handschriftliche Bemerkung:] vom Grad $64^{\mathrm{ti}}$.}. Mit indirekten Verfahren lassen sich jedoch die Werte der Unbekannten $x',x''$ recht \SourceLeaf{233}einfach aus diesen Gleichungen in ihrer vorliegenden Form gewinnen. Meist erhält man schon Näherungswerte der Unbekannten, wenn man zunächst $Q'$ und $Q''$ vernachlässigt. Es ergibt sich
\[
\begin{gathered}
x'=\frac{c''+d''(b''+c')+d'd''b'}{1-d'd''}\\[.6em]
x''=\frac{c'+d'(b'+c'')+d'd''b''}{1-d'd''}\mathpunct{\cdot}
\end{gathered}
\]
Sobald ein Näherungswert für eine der beiden Unbekannten vorliegt, lassen sich Werte, die den Gleichungen genau genügen, sehr leicht gewinnen. Sei $\xi'$ ein Näherungswert von $x'$, der beim Einsetzen in Gleichung I den Wert $x''=\xi''$ ergibt. Setzt man ebenso $x''=\xi''$ in Gleichung II ein, so ergibt sich $x'=X'$. Man wiederholt dieselben Rechnungen und setzt für $x'$ in I einen anderen Wert $\xi'+\nu'$ ein. Daraus ergibt sich $x''=\xi''+\nu''$, und setzt man diesen Wert in II ein, so erhält man $x'=X'+N'$. Der korrigierte Wert von $x'$ ist dann
\[
=\xi'+\frac{(\xi'-X')\nu'}{N'-\nu'}=\frac{\xi'N'-X'\nu'}{N'-\nu'}\mathpunct{\raisebox{.45ex}{,}}
\]
und der korrigierte Wert von $x''$
\[
=\xi''+\frac{(\xi'-X')\nu''}{N'-\nu'}\mathpunct{\cdot}
\]
Scheint es der Mühe wert, so wiederholt man dieselben Rechnungen mit dem korrigierten Wert von $x'$ und einem weiteren, davon leicht abweichenden Wert, bis sich Werte von $x',x''$ ergeben, die den Gleichungen I und II genau genügen. Im Übrigen wird jeder, der auch nur einigermaßen in der Analysis geübt ist, Mittel finden, die Rechnung abzukürzen.

Bei diesen Rechnungen lassen sich die irrationalen Größen $(x'x'+a'a')^{\frac32}$, $(x''x''+a''a'')^{\frac32}$ durch Einführung der Bögen $z',z''$ bequem berechnen. Ihre Tangenswerte sind der Reihe nach $\dfrac{a'}{x'}\mathpunct{\raisebox{.45ex}{,}}$ $\dfrac{a''}{x''}\mathpunct{\raisebox{.45ex}{,}}$ woraus folgt
\[
\begin{gathered}
\surd(x'x'+a'a'\mathpunct{\raisebox{-.5ex}{.}})=r'=\frac{a'}{\sin z'}=\frac{x'}{\cos z'}\\[.6em]
\surd(x''x''+a''a'')=r''=\frac{a''}{\sin z''}=\frac{x''}{\cos z''}\mathpunct{\cdot}
\end{gathered}
\]
Diese Hilfsbögen sind zwischen $0$ und $180^0$ zu wählen, damit $r',r''$ positiv werden. Sie stimmen offenbar mit den Bögen $C'B'$, $C''B''$ überein. Auf diese Weise erhält man daher sowohl $r'$ und $r''$ als auch die Lage der Punkte $C',C''$.

Diese Bestimmung der Größen $x',x''$ setzt voraus, dass\newline $a',a'',b',b'',c',c'',d',d'',Q',Q''$ bekannt sind. Die ersten vier dieser Größen ergeben sich aus den Angaben der Aufgabe, die nächsten vier hängen von $P',P''$ ab. Die \SourceLeaf{234}Größen $P',P'',Q',Q''$ lassen sich noch nicht genau bestimmen. Da aber gilt
\[
\tag*{III.}P'=\frac{t''-t'}{t'-t}\cdot\frac{(\eta01)}{(\eta12)}
\]
\[
\tag*{IV.}P''=\frac{t''-t'}{t'''-t''}\cdot\frac{(\eta23)}{(\eta12)}
\]
\[
\tag*{V.}\begin{aligned}
Q'={}&\tfrac12kk(t'-t)(t''-t')\cdot\frac{r'r'}{rr''}\cdot{}\\[.5em]
&\frac{1}{(\eta01)(\eta12)\cos\tfrac12(v'-v)\cos\tfrac12(v''-v)\cos\tfrac12(v''-v')}
\end{aligned}
\]
\[
\tag*{VI.}\begin{aligned}
Q''={}&\tfrac12kk(t''-t')(t'''-t'')\cdot\frac{r''r''}{r'r'''}\cdot{}\\[.5em]
&\frac{1}{(\eta12)(\eta23)\cos\tfrac12(v''-v')\cos\tfrac12(v'''-v')\cos\tfrac12(v'''-v'')}\mathpunct{\raisebox{.45ex}{,}}
\end{aligned}
\]
erhält man sofort die Näherungswerte
\[
\begin{gathered}
P'=\frac{t''-t'}{t'-t}\mathpunct{\raisebox{.45ex}{,}}\qquad P''=\frac{t''-t'}{t'''-t''}\\[.6em]
Q'=\tfrac12kk(t'-t)(t''-t'),\qquad Q''=\tfrac12kk(t''-t')(t'''-t''),
\end{gathered}
\]
auf denen die erste Rechnung aufbaut.
